\scanpage{3}
Structure

\unsure{4}$^\circ$) La donnée d'une sous-catégorie \inserted{\unsure{str.$^{\mathrm t}$} pleine} de $M$ (motifs effectifs), \unsure{stabl.} pr $\otimes$, pr sommes, \deleted{\unsure{quotients}} pr sous-\deleted{\unsure{tout}}objets, telle que \inserted{où $\Lambda(-1)$ soit effectif et que} pour $\forall E$, $\exists$ entier $n$ tel que $E(-n)\inserted{{}=E\otimes\Lambda(-1)^{\otimes n}}$ soit ``effectif''. Une telle structure sur $G_{\unsure{\mathbb{Q}}}$ en définit une \unsure{analogue} sur $G_{\mathbb{C}}$, et réciproqt; donc une telle notion de structure peut être considérée comme équivalente à la donnée d'un
\deleted{\unsure{partie de} $\unsure{K^+}(G)$ \unsure{bon cône}}
\deleted{cône $\unsure{X^+}$, $K^+(G)$}
\inserted{cône $K^+(G)$}
\deleted{\unsure{l'anneau} \unsure{des caractères} de $G_{\unsure{\mathbb{Q}}}$ $(\unsure{X^+}+\unsure{X^+}\subset\unsure{X^+})$}
\deleted{$\unsure{\Lambda}$ \unsure{stabl.} pr produit, tel que \unsure{pour tout}}
dans \inserted{l'anneau des car. rat/$\mathbb{Q}$,} $K(G)$, contenant $0$, \deleted{$\unsure{1}$} $\varepsilon$, stable pr multiplication, et \margin{\inserted{engendré pr les élém. de la base canonique\deleted{\unsure{s}} (classes de repr. irréd sur $\mathbb{Q}$) qu'il contient.}} À priori, on \unsure{aurait|devrait} là une notion de structure faisant intervenir la structure sur $\mathbb{Q}$.

Il y a lieu de poser d'autres axiomes. \inserted{Par exemple,} la condi. que pour $\forall$ \deleted{$\unsure{A}$} $E\in\mathrm{Ob}\,M^+$, $E$ soit à graduations \uuline{positives}.

\margin{De plus la condi. que \unsure{structure envisagée} provienne d'une structure de \deleted{$\unsure{=}$} type donné sur $G^0$, \unsure{stable par} \unsure{l'action de $G$ tout entier}.}

\margin{\uline{Axiome} Pour tt \deleted{\unsure{poids algébrique}} \inserted{\unsure{degré}} \unsure{$i\geq0$|il n'y a} \unsure{qu'un} nb fini de types de motifs effectifs de poids $i$ \ldots\ ???
En tous cas, on veut que les motifs effectifs de poids $0$ soient tels que $G^\circ$ opère en eux trivial$^{\mathrm t}$.}

Un \unsure{hom} $G_{m\mathbb{C}}\xrightarrow{i}G_{\mathbb{C}}$ \deleted{positif} \uline{effectif} si pour tout $E\in\mathrm{Ob}\,M^+$, la graduation correspondante de $E_{\mathbb{C}}$ est à

\mspage{2}
